\documentclass[11pt]{article}
\usepackage[margin=1in]{geometry}
\usepackage{amsmath,amssymb,bm,braket}
\usepackage{hyperref}
\begin{document}
\paragraph*{Problem setup:}
Consider a graphene layer sitting on a 3D substrate. The substrate has charged impurities of which the density is $n_i$ $\AA^{-3}$. Assume charge puddles with characteristic lateral screening length $\xi$ will form,  in which the 3D substrate and the sample layer together are charge neutral and form a thermal equilibrium. Such domains lead to electron and hole puddles. As a result, there will be a plateau in conductivity when changing the gate voltage $V_g$.

\paragraph*{Main problem:}

How will $\xi$ and the width of the plateau $\Delta V_g$ scale with the impurity density $n_i$?  Assuming $\xi\propto n_i^{\alpha}$, $\Delta V_g\propto n_i^{\beta}$, find $\alpha$ and $\beta$. Will such a plateau appear in a 3D topological insulator? Are the charged impurities still important there? Do charged impurities give long-range or short-range scattering? Will the long-range scattering in both graphene and a 3D topological insulator give longer mean free path than short-range scattering?



\paragraph*{Geometric and scaling assumptions}
The gate and graphene sheet have lateral dimensions much greater than $\xi$;
the substrate is thick on that scale. A puddle samples a substrate volume of
order $\xi^3$ and a graphene area of order $\xi^2$. The word ``domain'' denotes
this screening region, not a separately ordered thermodynamic phase.
Assume uncorrelated charged impurities, equilibrium local charge compensation,
and the self-consistent Thomas--Fermi screening regime of a two-dimensional
Dirac band. Hold gate capacitance per unit area, dielectric environment and
band parameters fixed when varying $n_i$. Treat $n_i$ as a volume density with
units of inverse length cubed. For the topological-insulator comparison,
consider its gapless surface Dirac states and nonmagnetic charged impurities.
\end{document}
